% !TeX root = ../main.tex
\section{The design for Guitar Virtual Amp}
\label{sec:design}

This section is the plan derived from Part~II, written before any code. Each
step carries its status today; Section~\ref{sec:impl} describes what is built
in detail.

\subsection{What the app already had}

\begin{itemize}
\item \textbf{A guitar stem per part}: BS-RoFormer (Demucs as fallback)
  separation~\cite{lu2023bsroformer,rouard2023htdemucs} and a split of the stereo
  guitar stem by pan position into parts that play different
  things~\cite{deangelo2026tabreview}.
\item \textbf{The notes and the grid}: onsets, offsets, pitches and techniques of
  every guitar part; beats, bars and tempo from Beat This!~\cite{foscarin2024beatthis}.
\item \textbf{An offline renderer of the real chain}: \code{backend::offline::render_at}
  runs the same pedal code as the live path (48\kHz, 16-bit samples, 256-sample
  blocks) at a chosen input RMS level. The full chain of 17 classic pedals costs
  about 430\,ns per sample on one core, a NAM WaveNet about 1.9\,\textmu s.
\item \textbf{Presets with song metadata}: 20 factory presets, user presets with
  \code{song} and \code{artist} fields, and a tone suggestion on the Tabs page
  that prefers a user preset tagged with the song (otherwise it guesses a factory
  preset from the track's instrument and name).
\item \textbf{Capture and cabinet players}: the NAM pedal and an IR convolver.
\end{itemize}

\subsection{Constraints}

\begin{itemize}
\item The chain is mono, works on 16-bit samples and has no gradients: the
  estimator treats it as a black box. Stereo effects cannot be reproduced.
\item The estimate may run offline on a desktop; its output, a preset, also
  works on Android.
\item The analysed audio must be the user's own: Spotify's developer policy
  forbids analysing its content, so the feature applies to the same local files
  and downloads as transcription.
\item Models that are distributed must be trained on data that allow it
  (Section~\ref{sec:datasets}); capture libraries can be searched by metadata,
  not bulk downloaded.
\item Patent CN118865922B claims a pipeline of similar shape~\cite{huanyin2025presetpatent};
  its claims should be reviewed before distributing the feature where the patent
  applies.
\end{itemize}

\subsection{The planned pipeline}

\begin{figure}[h]
\centering
\begin{tikzpicture}[
  node distance=2.4mm,
  box/.style={draw, rounded corners=2pt, font=\small, align=left,
              minimum height=7mm, inner sep=3pt, fill=white, text width=125mm},
  arr/.style={-{Stealth[length=2mm]}, semithick}]
  \node[box, fill=appc!10] (s1) {1. Separation, pan split, notes, beats \hfill \textsf{\scriptsize built (transcription)}};
  \node[box, below=of s1, fill=appc!10] (s2) {2. Pick an excerpt per part \hfill \textsf{\scriptsize built: loudest 12\,s}};
  \node[box, below=of s2, fill=appc!10] (s3) {3. Stand-in dry signal from the notes \hfill \textsf{\scriptsize built: pluck, calibrated on the player's guitar}};
  \node[box, below=of s3, fill=warnc!10] (s4) {4. Starting points: factory presets; learned guess \hfill \textsf{\scriptsize presets built; learned planned}};
  \node[box, below=of s4, fill=appc!10] (s5) {5. CMA-ES through the offline renderer, preset prior \hfill \textsf{\scriptsize built}};
  \node[box, below=of s5, fill=warnc!10] (s6) {6. Delay, reverb, modulation from notes and tempo \hfill \textsf{\scriptsize delay built}};
  \node[box, below=of s6, fill=appc!10] (s7) {7. Linear match: minimum-phase IR in a Cabinet \hfill \textsf{\scriptsize built}};
  \node[box, below=of s7, fill=appc!10] (s8) {8. Save as presets tagged with the song and track \hfill \textsf{\scriptsize built}};
  \foreach \a/\b in {s1/s2,s2/s3,s3/s4,s4/s5,s5/s6,s6/s7,s7/s8} \draw[arr] (\a) -- (\b);
\end{tikzpicture}
\caption{The planned tone-estimation pipeline and the status of each step.
Green: built and committed; orange: partly built or in progress.}
\label{fig:pipeline}
\end{figure}

The plan, step by step (Figure~\ref{fig:pipeline}):
\begin{enumerate}
\item \textbf{Separation and transcription} come from the tablature pipeline.
\item \textbf{Excerpts}: for each part, pick 5--20\,s where the part is loud
  relative to the rest of the mix, other guitars are quiet and the notes cover
  several registers; cluster excerpts by an effect embedding to find tone changes
  between sections and estimate one preset per cluster. \emph{Built so far}: the
  loudest 12\,s window.
\item \textbf{Stand-in dry signal}, from (a) a resynthesis of the transcribed
  notes, (b) an effect-removal model trained on the app's own renders, or (c) the
  player's own DI recorded while playing along in practice mode. (c) has the right
  guitar and the wrong timing; (a) the right timing and the wrong guitar; (b)
  both, with removal errors. \emph{Built}: (a). A fourth source, measured
  later, is a mix: (a) for the notes, with its colour calibrated on (c)
  (Section~\ref{sec:results-colour}).
\item \textbf{Starting points}: the closest factory preset, then a network
  trained on renders of the template chain, mixed over backing tracks and
  separated by the same BS-RoFormer so it sees separation artefacts (the StemFX
  recipe~\cite{cheng2026stemfx}), sharing each dry clip across many chains so it
  learns the tone rather than the riff~\cite{nathanben2026tonerunner}. AFx-Rep
  (Apache-2.0)~\cite{steinmetz2024stito} is a starting point for its encoder.
\item \textbf{Refinement} by CMA-ES. The plan lists four terms for the
  objective: an aligned spectral term, an alignment-free term (long-term spectra
  per register, crest factor, compression statistics), a small-weight
  effect-embedding term, and a prior over factory and user
  presets~\cite{yu2025improving}. \emph{Built}: the aligned spectral term, with an
  EQ allowance, and a prior over the factory presets.
\item \textbf{Time-based effects from the notes}: delay time from echoes of
  onsets (or candidate values tied to the tempo), reverb decay from rests after a
  note ends, modulation rate from sustained notes or an LFO
  extractor~\cite{mitcheltree2023modulation}. Measured values fix or initialise
  the search, which avoids rugged landscapes. \emph{Built}: delay.
\item \textbf{Linear match}: the smoothed difference between the long-term
  spectra of the stem and of the rendered chain, as a minimum-phase IR in a
  Cabinet pedal with a gain limit, as Tone Match does~\cite{fractalndtonematch}.
\item \textbf{Output}: a user preset with song and artist set, so the existing
  tone suggestion offers it; other candidates kept; settings recorded for later
  evaluation.
\end{enumerate}
\textbf{Captures} (later): NAM files on the user's disk as candidates for the amp
slot, ranked by rendering the dry signal through each at several input
gains~\cite{sd4562026nameqmatch}; TONE3000 queried by metadata for candidates the
user then chooses to download. Learning a capture from the stem alone by unpaired
or diffusion-based fitting~\cite{wright2023adversarial,moliner2025unsupervised}
is research, at hours of GPU time per tone.

\subsection{Why this design}

\begin{itemize}
\item \textbf{Against prediction only}: synthetic-to-real gaps and ambiguous
  settings; refinement by rendering gained 4.9\dB{} SI-SDR in Okita and
  Katayose~\cite{okita2026audio}.
\item \textbf{Against an embedding-only search}: ST-ITO's authors report failure
  on guitar tone, and TAPE shows the objective drifting~\cite{steinmetz2024stito,long2026tape};
  a dry signal aligned with the target allows a spectral objective instead.
\item \textbf{Against neural tone transfer}: EG-VAE gives the best guitar tone
  similarity but returns audio, has no released code, and cannot run in the
  real-time chain~\cite{yeh2026egvae}.
\item \textbf{Against linear matching only}: a filter cannot set gain or
  dynamics, and Tone Match's own manual requires the gain to be right
  first~\cite{fractalndtonematch}.
\item \textbf{For the app's own chain}: its parameters are what the player
  edits, its renderer is the exact forward model, and training on its renders
  gives labels in its own parameter space with licences under our control.
\item \textbf{For a dry signal made from the notes}: the transcription already
  exists, and it gives the one thing the other stand-ins lack, alignment with the
  target. No published work does this.
\end{itemize}
